%% ============================================================
%%  Chapter 2 — Background
%% ============================================================
\chapter{Background}
\label{ch:background}

\section{Kubernetes and Cloud-Native Networking}
\label{sec:bg_kubernetes}

\subsection{Kubernetes Architecture}
Kubernetes~\cite{kubernetes2014} is an open-source container orchestration platform
that automates the deployment, scaling, and lifecycle management of containerised
workloads. A Kubernetes cluster comprises a \emph{control plane} (responsible for
scheduling, state management, and API serving) and a set of \emph{worker nodes} that
execute application workloads encapsulated in \emph{Pods}.

Network communication in Kubernetes follows a flat networking model: every Pod is
assigned a unique IP address and can communicate directly with any other Pod across nodes
without NAT, via a \emph{Container Network Interface (CNI)} plugin. CNI plugins such as
Flannel, Calico, Weave, and \emph{Cilium}~\cite{cilium} implement this networking model
at the OS level.

\subsection{Kubernetes Networking Constructs}
\begin{description}
  \item[Pod:] The atomic deployment unit; one or more co-located containers sharing a
  network namespace.
  \item[Service:] A stable virtual IP (\texttt{ClusterIP}) that load-balances traffic
  across a set of pods.
  \item[Namespace:] A logical isolation boundary grouping related workloads.
  \item[NetworkPolicy:] A Kubernetes-native API object that specifies ingress/egress
  rules for pods. Enforcement is CNI-dependent.
\end{description}

\subsection{Cilium and Hubble}
Cilium~\cite{cilium} is a CNI plugin that uses eBPF to implement networking, security,
and observability entirely in the Linux kernel. Unlike iptables-based CNIs, Cilium
attaches eBPF programs to network interfaces, enabling high-performance packet processing
without context switches to userspace.

Hubble~\cite{hubble} is Cilium's distributed observability layer. It captures
\emph{flow records} --- structured metadata about every network connection (source pod,
destination pod, protocol, bytes, packets, verdict, latency) --- from the eBPF data plane
and exposes them via a gRPC API and a web UI. In this project, Hubble serves as the
primary telemetry source for graph construction.

% ─────────────────────────────────────────────
\section{eBPF: Extended Berkeley Packet Filter}
\label{sec:bg_ebpf}

\subsection{Historical Context}
The Berkeley Packet Filter (BPF) was introduced by McCanne and Jacobson in 1992~\cite{mccanne1993bpf}
as an in-kernel packet filtering mechanism for \texttt{libpcap}/\texttt{tcpdump}. It
defined a simple register-based virtual machine running within the kernel to evaluate
filter expressions on packet data without copying packets to userspace.

Extended BPF (eBPF), introduced in Linux 3.15 (2014) and stabilised through 4.x and 5.x
kernels, fundamentally transformed BPF from a packet filter into a general-purpose,
safe, sandboxed programmability layer for the Linux kernel~\cite{gregg2019bpf}.

\subsection{eBPF Architecture}
\label{subsec:ebpf_arch}

The eBPF subsystem comprises four core components:

\begin{enumerate}
  \item \textbf{eBPF Programs:} Written in a restricted subset of C and compiled to
  eBPF bytecode using LLVM/Clang. Program types determine the kernel hook they attach to
  (e.g., \texttt{BPF\_PROG\_TYPE\_XDP}, \texttt{BPF\_PROG\_TYPE\_TRACEPOINT},
  \texttt{BPF\_PROG\_TYPE\_LSM}).

  \item \textbf{Verifier:} A static analyser that exhaustively proves safety properties
  before a program is loaded into the kernel: no unbounded loops, no invalid memory
  accesses, no out-of-bounds reads, and type safety of all helper calls.

  \item \textbf{JIT Compiler:} Translates verified eBPF bytecode to native machine code
  (x86-64, ARM64, etc.) for near-native execution performance.

  \item \textbf{BPF Maps:} Kernel-resident key-value stores that provide persistent state
  and bidirectional communication between eBPF programs and userspace applications.
  Common map types include hash maps, LRU maps, ring buffers, and per-CPU arrays.
\end{enumerate}

\subsection{Key eBPF Hook Points}
\label{subsec:ebpf_hooks}

\begin{table}[h]
\centering
\caption{Common eBPF program types and their kernel attachment points}
\label{tab:ebpf_hooks}
\begin{tabularx}{\textwidth}{lXl}
\toprule
\textbf{Program Type} & \textbf{Attachment Point / Use Case} & \textbf{Typical Latency} \\
\midrule
XDP & Driver/NIC level; earliest ingress hook; line-rate packet processing & $<$1~$\mu$s \\
TC (Traffic Control) & \texttt{cls\_bpf} at TC ingress/egress; bidirectional & $\sim$1--5~$\mu$s \\
Socket Filter & Per-socket packet filtering (e.g., \texttt{AF\_PACKET}) & Low \\
kprobe/kretprobe & Dynamic kernel function instrumentation & Moderate \\
Tracepoint & Static kernel tracepoints (stable ABI) & Low \\
LSM & Linux Security Module hooks; access control enforcement & Low \\
Cgroup & Per-cgroup network/socket control & Low \\
\bottomrule
\end{tabularx}
\end{table}

\subsection{eXpress Data Path (XDP)}
XDP~\cite{hoiland2018xdp} is an eBPF program type that attaches at the earliest possible
point in the packet receive path --- directly in the NIC driver or at the TC ingress hook.
XDP programs return an action verdict for each packet:
\begin{itemize}
  \item \texttt{XDP\_DROP} — discard the packet immediately.
  \item \texttt{XDP\_PASS} — forward to the kernel network stack.
  \item \texttt{XDP\_TX} — retransmit the packet out the same interface.
  \item \texttt{XDP\_REDIRECT} — redirect to another interface or CPU.
\end{itemize}
XDP can process millions of packets per second on commodity hardware, making it suitable
for high-rate DDoS mitigation~\cite{hoiland2018xdp}.

\subsection{eBPF Safety Guarantees}
A critical property that distinguishes eBPF from kernel modules is its safety model: all
programs must pass verifier analysis before execution. The verifier proves termination
(no loops without bounded iteration), memory safety, and valid register type propagation.
This allows untrusted eBPF programs to be loaded without risking kernel stability.

% ─────────────────────────────────────────────
\section{DDoS Attacks: Taxonomy and Characteristics}
\label{sec:bg_ddos}

\subsection{Taxonomy}
DDoS attacks can be classified along several dimensions~\cite{ddos_survey_2022}:

\begin{description}
  \item[By OSI Layer:]
    \textit{Network-layer} attacks target Layers 3--4 (e.g., ICMP flood, SYN flood,
    UDP amplification). \textit{Application-layer} attacks target Layer 7 (e.g., HTTP
    flood, Slowloris, DNS query flood).

  \item[By Volume vs. Sophistication:]
    \textit{Volumetric} attacks saturate bandwidth or connection tables with high-rate
    traffic. \textit{Low-and-slow} attacks (e.g., Slowloris~\cite{slowloris}) send traffic
    at rates indistinguishable from legitimate users but exhaust server connection state
    over time.

  \item[By Source:]
    \textit{Botnet-driven} attacks coordinate thousands of compromised hosts.
    \textit{Amplification} attacks exploit protocols with large response-to-request ratios
    (DNS, NTP, SSDP) using spoofed source IPs.
\end{description}

\subsection{DDoS in Kubernetes Environments}
Kubernetes-specific threat vectors include~\cite{k8s_security_2023}:
\begin{itemize}
  \item Compromised pods acting as internal attack sources (East-West DDoS).
  \item Exploiting Kubernetes auto-scaling to amplify attack cost for defenders.
  \item Targeting specific microservices that act as bottlenecks (e.g., authentication
  services, databases).
  \item Leveraging service mesh encryption to hide attack traffic from perimeter sensors.
\end{itemize}

% ─────────────────────────────────────────────
\section{Graph Theory Foundations}
\label{sec:bg_graphs}

\subsection{Graph Definitions}
A \emph{graph} $G = (V, E)$ consists of a set of vertices $V$ and edges $E \subseteq
V \times V$. In this project, a \emph{communication graph} $G_t$ is constructed at time
$t$ where:
\begin{align}
  V &= \{\text{pods, services, external endpoints}\} \\
  E &= \{\text{network flows observed in window } [t-\Delta t, t]\}
\end{align}
Node features $\mathbf{x}_v \in \mathbb{R}^d$ encode per-pod metrics (request rate, error
rate, CPU utilisation). Edge features $\mathbf{e}_{uv} \in \mathbb{R}^k$ encode flow
metrics (throughput, latency, packet count).

\subsection{Dynamic Graphs}
Real Kubernetes workloads exhibit non-stationary communication patterns: pods are created
and destroyed, traffic patterns shift with user load, and deployments change topology.
A \emph{dynamic graph} or \emph{temporal graph} $\{G_{t_1}, G_{t_2}, \ldots, G_{t_n}\}$
captures this temporal evolution~\cite{kazemi2020representation}.

% ─────────────────────────────────────────────
\section{Graph Neural Networks}
\label{sec:bg_gnns}

\subsection{Message Passing Framework}
The dominant paradigm for GNNs is \emph{message passing}~\cite{gilmer2017mpnn}. Each
node $v$ iteratively aggregates information from its neighbourhood $\mathcal{N}(v)$:
\begin{equation}
  \mathbf{h}_v^{(l+1)} = \text{UPDATE}\left(\mathbf{h}_v^{(l)},
    \text{AGGREGATE}\left(\{\mathbf{h}_u^{(l)} : u \in \mathcal{N}(v)\}\right)\right)
\end{equation}
where $\mathbf{h}_v^{(l)}$ is the embedding of node $v$ at layer $l$. After $L$ layers,
$\mathbf{h}_v^{(L)}$ encodes information from the $L$-hop neighbourhood of $v$.

\subsection{Graph Convolutional Network (GCN)}
GCN~\cite{kipf2017gcn} instantiates the message passing framework with:
\begin{equation}
  \mathbf{H}^{(l+1)} = \sigma\left(\tilde{\mathbf{D}}^{-\frac{1}{2}}
    \tilde{\mathbf{A}} \tilde{\mathbf{D}}^{-\frac{1}{2}} \mathbf{H}^{(l)} \mathbf{W}^{(l)}\right)
\end{equation}
where $\tilde{\mathbf{A}} = \mathbf{A} + \mathbf{I}$ is the adjacency matrix with
self-loops, $\tilde{\mathbf{D}}$ is the degree matrix, and $\mathbf{W}^{(l)}$ is a
learnable weight matrix.

\subsection{GraphSAGE}
GraphSAGE~\cite{hamilton2017graphsage} uses inductive learning by sampling and
aggregating fixed-size neighbourhoods, enabling generalisation to unseen nodes --- critical
for Kubernetes environments where pods are created dynamically.

\subsection{Graph Autoencoders (GAE)}
GAE~\cite{kipf2016vgae} comprises an encoder $q(Z|X,A)$ that maps node features to a
latent space $Z$, and a decoder $p(A|Z)$ that reconstructs the adjacency matrix. Anomaly
detection uses the reconstruction error:
\begin{equation}
  \text{score}(v) = \left\| \mathbf{A}_v - \hat{\mathbf{A}}_v \right\|_F
\end{equation}
High scores indicate structural anomalies --- edges that should not exist (attack traffic
from unexpected sources) or edges that disappear (service degradation).

\subsection{Temporal Graph Networks (TGN)}
TGN~\cite{rossi2020tgn} processes a stream of timestamped interaction events, maintaining
per-node \emph{memory states} that capture longitudinal behaviour. This enables detection
of attacks whose signatures emerge only over extended time windows.

% ─────────────────────────────────────────────
\section{Summary}
This chapter has established the foundational concepts required to understand the
remainder of this report: Kubernetes networking, the eBPF subsystem and its hook model,
the DDoS threat landscape in cloud-native environments, and the graph ML techniques that
will be applied for anomaly detection. Chapter~\ref{ch:literature} surveys how these
technologies have been applied in prior research and identifies the gaps that motivate
this project.
